\section{A sharper bound for exact counting}\label{app:quantitative}

We prove Theorem~\ref{thm:main} by contracting Boolean tensor networks.
The two running times reflect the number of wires in the networks.
A formula with $m\le Dn$ clauses gives a network with $O_{k,D}(n)$
wires. For arbitrary formulas, we use the parity tests of
Section~\ref{sec:hashes}. The degree bound of Lemma~\ref{lem:total-order}
lets us index a monomial retained in the product of their polynomial
lifts with $O_k(n\log\log n)$ wires. We give both encodings before
proving the contraction theorem they use.

\subsection{Boolean tensor contraction}
\label{subsec:tensor-interface}

A Boolean tensor network is a bipartite multigraph $G$ in which each
edge carries a single bit. Each vertex $v$ has an ordered list
$\delta(v)$ of incident edges and a tensor $T_v$ assigning a ring
element to each Boolean assignment to those edges. Parallel edges
have distinct ports. The \emph{contraction} of the network is
\[
 Z=\sum_{x\in\{0,1\}^{E(G)}}\prod_{v\in V(G)} T_v(x|_{\delta(v)}).
\]
An isolated vertex contributes its scalar entry. We measure the size
of the network by its number of edges.

\begin{theorem}[Boolean tensor contraction]\label{thm:tensor-contraction}
Fix positive constants $C,A,b$.  Let $s\ge2$ be an integer, and let
$\mathcal L$ be the total length of the graph and tensor descriptions
of a bipartite tensor network with at most $Cs$ edges and maximum
degree at most $\lceil s^{9/10}\rceil$.  Suppose that the number of
vertices is bounded by a fixed polynomial in $s$, and that the entries
belong to
\[
 \mathcal B_{P_0,\ell}
   =(\mathbb Z/2^{P_0}\mathbb Z)[V]/(V^\ell),
 \qquad 1\le P_0,\ell\le Cs.
\]
Suppose, uniformly for every vertex and Boolean choice of its ports,
that its entry can be computed from the description in time
$(s+\mathcal L)^b2^{A s^{9/10}}$ and space $(s+\mathcal L)^b$.
Then the contraction can be computed deterministically in
\[
 (s+\mathcal L)^{O(1)}
  2^{O(s/\sqrt{\log s})}
\]
bit operations and $(s+\mathcal L)^{O(1)}$ bits of space.  Constants
depend only on the fixed bounds in the hypotheses.  In particular,
$\ell=1$ is permitted.
\end{theorem}

The algorithm in Section~\ref{app:tensor} replaces a single global
Taylor expansion by recursion on the wires. At each call it deforms
the wire identifications so that, at a chosen interpolation node,
only ports in the same group interact. Near that node, interactions
between groups expand in variables called modes. A mode of weight
$j$ has coefficient divisible by $2^j$, so only monomials of total
weight below the working precision can contribute. Tapering the
Hermite multiplicities with the two-adic valuation of the node
bounds that precision. With $M$ wires and $u$ retained valuation
levels, the number of mode monomials is $2^{O(M/\sqrt{u})}$, and
each monomial factors into contractions with at most $M/8$ wires.
At recursion depth $d$, $M=O(s8^{-d})$ and $u=\Theta((\log s)/2^d)$.
The exponents therefore sum geometrically to $O(s/\sqrt{\log s})$.

\subsection{Fixed clause density}\label{subsec:sharp-fixed-density}

\begin{proof}[Proof of the fixed-density bound in Theorem~\ref{thm:main}]
Give each clause a vertex carrying its Boolean satisfaction table, with
one port for each distinct variable in the clause. A tautology may be
dropped, and an empty clause gives count zero. For a variable with
$d\ge1$ ports, require their bits to agree. For $d=1$ attach the unary
tensor $(1,1)$, and for $d=2$ attach a binary equality tensor. For
$d\ge3$ take a tree with $d$ external ports and trivalent equality
tensors at its internal vertices. Such a tree has $d-2$ internal
vertices, and each common value on its ports extends uniquely to all
its internal edges.

There are at most $km$ clause ports. The graph has $O_k(n+m)$ edges
and vertices, and maximum degree at most $\max\{k,3\}$. Subdivide
every edge once by a binary equality tensor. The old vertices and the
new vertices form the two sides of a bipartite network. The two bits
on a subdivided edge must be equal, so this preserves the contraction.
Every satisfying assignment on the occurring variables contributes
one and every other assignment contributes zero.

With $s=n+m$, apply Theorem~\ref{thm:tensor-contraction} over
$\mathbb Z/2^{n+1}\mathbb Z$, taking $P_0=n+1$ and $\ell=1$. The
local tables have bounded size for fixed $k$. The
degree and precision bounds hold for all sufficiently large $s$, and
bounded instances can be enumerated directly. If $u$ declared
variables are unused, multiply the contraction by $2^u$ modulo
$2^{n+1}$. Its least nonnegative residue is $\#F$, since
$0\le\#F\le2^n$. This includes the empty formula, for which the
empty network has contraction one. The time is
$s^{O_k(1)}2^{O_k(s/\sqrt{\log s})}$ and the space is polynomial.
When $m\le Dn$, we have $n\le s\le(1+D)n$, giving exponent
$O_{k,D}(n/\sqrt{\log n})$.
\end{proof}

\subsection{Arbitrary clause density}\label{subsec:sharp-dense}

The sparsification and greedy row selection in Section~\ref{sec:hashes}
reduce this case to counting common zeros of sparse polynomial equations.
We retain the parity lift of Section~\ref{sec:setup}. Its degree bound
lets us encode a monomial by a subset of test--variable incidences using
$O(n\log\log n)$ Boolean wires.

\begin{lemma}[Counting sparse polynomial systems by contraction]
\label{lem:tensor-sparse-query}
Fix integers $k,c\ge1$. Let $P_1,\ldots,P_a$ be polynomials over
$\mathbb F_2$ on $n\ge20$ declared Boolean variables. Suppose
$a\le cn$ and each $P_i$ is listed as at most
$\lceil\log_2 n\rceil^c$ monomials of degree at most $k$. The number
of common zeros is computable deterministically in
$n^{O_{k,c}(1)}2^{O_{k,c}(n\log\log n/\sqrt{\log n})}$ bit operations and
$n^{O_{k,c}(1)}$ bits of space.
\end{lemma}

\begin{proof}
Put $\lambda=\lceil\log_2 n\rceil$ and
$\nu=\lceil\log_2\lambda\rceil$.
Within each monomial replace positive powers of a variable by its
first power, and combine duplicate monomials modulo two. The scope
of test $i$ consists of the variables in its remaining monomials.
It has size at most $k\lambda^c$.
Let $N_0$ be the sum of the scope sizes, counting a pair $(i,v)$
once if variable $v$ belongs to test $i$. Thus $N_0\le ak\lambda^c$.
Fix, in all possible ways, the variables belonging to more than
\[
 D_*=\max\{1,\lceil N_0/n\rceil\}\lceil\sqrt{\lambda}\rceil
\]
scopes. There are at most $n/\sqrt{\lambda}$ such variables. Their
assignments partition the original count into at most
$2^{n/\sqrt{\lambda}}$ branch counts. On each branch substitute the
fixed bits and combine duplicate monomials again. The free scope
$I_i$ of test $i$ still has size at most $k\lambda^c$, and each free
variable appears in at most $D_*=\lambda^{O_{k,c}(1)}$ scopes. Free
variables absent from every scope remain declared. A constant-one
equation makes the branch count zero. We may discard constant-zero
equations, and if none remain the answer for the branch is two to the
power of its number of free variables. Below $a$ denotes the number
of remaining equations and $n$ the original variable count.

\paragraph{The polynomial on the incidences.}
Give each pair $(i,v)$ with $v\in I_i$ its own indeterminate
$z_{iv}$. Write $z_i=(z_{iv})_{v\in I_i}$. Form
$\widetilde L_i(z_i)$ by summing the remaining monomials of the
normalized branch polynomial with integer coefficients, using $z_{iv}$
for $x_v$. We use the ring
\[
 \mathcal B=(\mathbb Z/2^p\mathbb Z)[V]/(V^\ell),
 \qquad p=\ell=n+1 .
\]
Take $h(Y)=Y(Y-1)$ and the integer polynomials $e_j$ from
Lemma~\ref{lem:rational-idempotent}. Reduce their coefficients
modulo $2^p$, and form the factor
\[
 B_i(z_i,V)=\sum_{j=0}^{\ell-1}
   e_j(\widetilde L_i(z_i))\,h(\widetilde L_i(z_i))^j V^j .
\]
Let $\overline B_i$ be its multilinear representative in the $z_i$,
obtained in each coefficient of $V$ by replacing each positive
power of an incidence indeterminate by its first power. It agrees
with $B_i$ on every Boolean assignment to $z_i$. Define
$F(z,V)=\prod_i\overline B_i(z_i,V)$ in $\mathcal B[z]$. The
incidence indeterminates for different tests are disjoint, so $F$
is multilinear even before substituting the same Boolean $x_v$ for
all copies of variable $v$.

The coefficient of $V^j$ in $\overline B_i$ has incidence degree
at most $k(1+2j)$, since $\deg e_j\le1$, $\deg h=2$, and
multilinearization only lowers monomial degrees. Multiplication in
$\mathcal B$ retains only terms with $\sum_i j_i\le\ell-1$.
Consequently every incidence monomial in $F$ has degree at most
\begin{equation}\label{eq:sharp-incidence-degree}
 K=k\bigl(a+2(\ell-1)\bigr)=O_{k,c}(n).
\end{equation}
Let $\mathcal S_\ell$ sum the coefficients of $1,V,\ldots,V^{\ell-1}$
in $\mathcal B$. At $z_{iv}=x_v$, the multilinearized factor
$\overline B_i$ agrees with $B_i$ for the normalized branch polynomial.
The integer $\widetilde L_i((x_v)_{v\in I_i})$ has the same parity
as the original $P_i$ evaluated at $x$ and the fixed bits. Applying
Lemma~\ref{lem:total-order} to these integers shows that the branch
count is congruent modulo $2^p$ to
\begin{equation}\label{eq:sharp-branch-residue}
 \mathcal S_\ell\left(\sum_x F\bigl((z_{iv}=x_v)_{iv},V\bigr)\right),
\end{equation}
where $x$ runs over all free Boolean assignments. We apply the
coefficient-sum map after multiplying the test factors in the
truncated ring.

\paragraph{Indexing the monomials.}
Partition the tests and the free variables separately into
$J=\lceil n^{1/4}\rceil$ groups of sizes differing by at most one
on each side. Write $\mathcal T_g$ for a test group and
$\mathcal X_h$ for a variable group. A group contains
$O_c(n^{3/4})$ objects, so the number of its incidences is at most
$n^{3/4}\lambda^{O_{k,c}(1)}$ on either side. Let $N_{gh}$ be the
number of incidences from $\mathcal T_g$ to $\mathcal X_h$ and let
$N=\sum_{g,h}N_{gh}\le N_0$.

For an incidence monomial let $d_{gh}$ be the number of its
selected incidences joining groups $g,h$. Enumerate arrays of
integers $(d_{gh})$ with $0\le d_{gh}\le N_{gh}$ and
$\sum_{g,h}d_{gh}\le K$. There are at most
\[
 (K+1)^{J^2}=2^{O_{k,c}(\sqrt n\log n)}
\]
arrays, which can be enumerated using $J^2$ counters. For a fixed
array there are
\[
 C_{gh}=\binom{N_{gh}}{d_{gh}}
\]
subsets of the specified size for this pair. Draw
$\lceil\log_2 C_{gh}\rceil$ Boolean wires between the two group
vertices to index them in lexicographic order. For $C_{gh}=1$
both endpoints use the unique subset (possibly the entire incidence
set) with no wires. Any unused bit string gives tensor entry zero.
To decode a valid index, take the incidences in order and count the
possible completions with or without each incidence using binomial
coefficients. These have polynomial bit length, so decoding uses
polynomial time and space.

The union of the subsets over all pairs specifies a squarefree
monomial of degree at most $K$ in $N$ indeterminates. Including
all monomials of this degree bound gives the upper bound
$\binom{N+K}{K}$ by stars and bars. Thus the number
of wires $M$ in the network for this array satisfies
\begin{align}
 M &\le J^2+\log_2\prod_{g,h} C_{gh}
      \le J^2+\log_2\binom{N+K}{K} \notag\\
   &\le J^2+K\log_2\!\left(\frac{\mathrm e(N+K)}{K}\right)
      = O_{k,c}(n\nu).
 \label{eq:sharp-wire-budget}
\end{align}
Here $K=\Theta_{k,c}(n)$ and $N\le n\lambda^{O_{k,c}(1)}$.
Also $\lceil\log_2 C_{gh}\rceil\le N_{gh}$, so the degree of each
vertex is at most its number of incidences,
$n^{3/4}\lambda^{O_{k,c}(1)}$.

\paragraph{Tensor entries and contraction.}
A valid assignment of wire bits determines a subset $S_i\subseteq I_i$
at each test. Write $z_{i,S_i}=\prod_{v\in S_i}z_{iv}$. The entry
at a test group $g$ is the coefficient of its indicated monomial
in $\prod_{i\in\mathcal T_g}\overline B_i(z_i,V)$. Since the $z_i$
are disjoint, this is
\[
 \prod_{i\in\mathcal T_g}[z_{i,S_i}]\overline B_i(z_i,V)
 \qquad\text{in }\mathcal B.
\]
For a variable group $h$, put $b_v=|\{i:v\in S_i\}|$
for $v\in\mathcal X_h$. Its entry is
\[
 \sum_{x\in\{0,1\}^{\mathcal X_h}}
        \prod_{v\in\mathcal X_h}x_v^{b_v}
   =2^{|\{v\in\mathcal X_h:b_v=0\}|}\pmod{2^p}.
\]
Indeed a positive exponent requires $x_v=1$ and a zero exponent
leaves $x_v$ free. This includes variables outside all test scopes.

Each monomial of $F$ occurs in exactly one array and has exactly
one valid index at each group pair. Its coefficient in $F$ is the
product of the test-group coefficients. After the incidences with
the same $v$ are set equal, its sum over Boolean assignments is
the product of the variable-group entries. By
\eqref{eq:sharp-incidence-degree}, summing over these arrays exhausts
all monomials. Consequently the element of $\mathcal B$ inside
$\mathcal S_\ell$ in \eqref{eq:sharp-branch-residue} is the sum of
the contractions of these bipartite networks.

To evaluate one test-side coefficient, let $\mathbf1_T$ be the
Boolean vector on $I_i$ with support $T\subseteq S_i$. At
$\mathbf1_T$, the multilinear polynomial evaluates to the sum of its
coefficients indexed by subsets of $T$, and it agrees there with
$B_i$. Inverting this subset sum gives
\[
 [z_{i,S_i}]\overline B_i(z_i,V)=
 \sum_{T\subseteq S_i}(-1)^{|S_i|-|T|}B_i(\mathbf1_T,V)
 \qquad\text{in }\mathcal B.
\]
Each term costs polynomially many bit operations. Evaluate
$\widetilde L_i(\mathbf1_T)$ from the listed monomials and use the
polynomials $e_j$ of Lemma~\ref{lem:rational-idempotent}, which
are constructible in polynomial time. Store the $\ell$ coefficients
of $B_i(\mathbf1_T,V)$ modulo $2^p$. There are at most
$2^{|S_i|}\le 2^{k\lambda^c}$ subsets to process. Sum one test
at a time and multiply the resulting coefficients in $\mathcal B$
over the group. This evaluates a test-group entry in
$2^{\lambda^{O_{k,c}(1)}}$ bit operations and polynomial space.
The variable-side entries and checks for invalid wire indices take
polynomial time. The graph, degree array, and entry procedures have
polynomial-size descriptions.

Apply Theorem~\ref{thm:tensor-contraction} with size parameter
$s=n\nu$, root precision $p=n+1$, and truncation order $\ell=n+1$.
The graph has $2J$ vertices and $O_{k,c}(s)$ wires by
\eqref{eq:sharp-wire-budget}. For all sufficiently large $n$ its
maximum degree is at most $\lceil s^{9/10}\rceil$, its precision
and truncation order are $O(s)$, and its entry-evaluation time is
at most $2^{O_{k,c}(s^{9/10})}$. Each contraction then takes
$n^{O_{k,c}(1)}2^{O_{k,c}(s/\sqrt{\log s})}$ bit operations and
polynomial space. Here $s/\sqrt{\log s}=O(n\nu/\sqrt\lambda)$.
The high-incidence branches and degree arrays multiply the
running time by at most
\[
 2^{\,n/\sqrt{\lambda}+O_{k,c}(\sqrt n\log n)}.
\]
Both terms in this exponent are $O_{k,c}(n\nu/\sqrt{\lambda})$.
Process the branches and arrays sequentially. On each branch,
sum the contractions in $\mathcal B$ and apply $\mathcal S_\ell$.
Adding the branch residues modulo $2^{n+1}$ gives the original
count, since it lies between $0$ and $2^n$. This proves the
claimed bounds, using direct enumeration for bounded exceptional
sizes.
\end{proof}

\begin{proof}[Proof of the arbitrary-density bound in
Theorem~\ref{thm:main}]
Normalize the formula on its original $n$ declared variables as
in Section~\ref{subsec:general-reduction}. Empty clauses and
formulas with no clauses have immediate answers. For $n\ge20$
apply Lemma~\ref{lem:sparsify} with
$\sigma=1/\lceil\sqrt{\lambda}\rceil$, where
$\lambda=\lceil\log_2 n\rceil$. The tree has
$n^{O_k(1)}2^{O(n/\sqrt{\lambda})}$ nodes and polynomial depth.
Every leaf has occurrence bound
$B=O_k((\sigma^{-1}\log(2/\sigma))^{k-1})=\lambda^{O_k(1)}$
and at most $Bn$ nonempty clauses.

For a nontrivial leaf with $m_G$ clauses, use the initial parity
checks of Lemma~\ref{lem:sparse-checks} and the greedy procedure
of Lemma~\ref{lem:greedy-hashes}. Each initial row has support size
at most $D_0$, and each additional row has list length $P$.
Both bounds are $\lambda^{O_k(1)}$. By \eqref{eq:hash-query}, every row
gives a degree-at-most-$k$ equation with at most
$2^k\max\{P,D_0\}$ monomials, and every queried system has at
most $6n$ equations. A single constant $c=c(k)$ bounds both the
ratio of equations to variables and the exponent of $\lambda$ in
the monomial bound for all sufficiently large $n$.
At each greedy step, use Lemma~\ref{lem:tensor-sparse-query} to
compute $N(A;h_\gamma)$ for every candidate row $h_\gamma$ and
append one attaining the minimum. The evaluator applies also to
discarded candidates whose systems still accept nonsatisfying
assignments. One final call gives $N(A)$ for the selected matrix,
which equals the leaf count by Lemma~\ref{lem:greedy-hashes}.

Constructing the initial checks and enumerating candidates costs
$2^{\lambda^{O_k(1)}}$ bit operations apart from the evaluations.
There are at most $3n(2m_G)^P+1=2^{\lambda^{O_k(1)}}$ evaluations
per leaf. Since every fixed power of $\lambda$ is
$o(n\log\log n/\sqrt{\log n})$, these costs and the number of leaves
are absorbed in the claimed time bound
\[
 (n+m)^{O_k(1)}2^{O_k(n\log\log n/\sqrt{\log n})}.
\]
Here $m$ counts the input clauses. Traverse the tree
depth first and process one query at a time. The leaf solution sets
are disjoint, so their exact counts add with an $(n+1)$-bit
accumulator. Together with polynomial space per query and per
tree path, this gives the asserted polynomial-space bound.
\end{proof}

A single algorithm attains both bounds. With
$\lambda=\lceil\log_2 n\rceil$ and $\nu=\lceil\log_2\lambda\rceil$,
use the clause network when $m\le n\nu$ and the parity queries
otherwise. In the first branch $n\le s=n+m=O(n\log\log n)$,
so its exponent is $O_k(n\log\log n/\sqrt{\log n})$ even at growing
density. For each fixed $D$ and all sufficiently large $n$, every
input with $m\le Dn$ takes that branch with $s=O_D(n)$.
